% Part: counterfactuals
% Chapter: minimal-change-semantics
% Section: transitivity

\documentclass[../../../include/open-logic-section]{subfiles}

\begin{document}

\olfileid{cnt}{min}{tra}

\olsection{انتقاليت}

د مادي شرطي لپاره د زنځير قاعده برقرار ده: $!A \lif !B, !B
\lif !C \Entails !A \lif !C$۔ په بله وينا، مادي شرطي انتقالي
ده۔ د واقعيت خلاف شرطونو لپاره د سټالنکر په بېلګه کښې دوه مقدمې او يوه پايله لرو۔
\begin{explain}
  لومړۍ مقدمه د هوور د زېږېدو ځاے بدلوي۔ په دې فرضي حالت
  کښې د شوروي اتحاد په چاپېريال کښې د هغه کمونست کېدل منو۔

  دويمه مقدمه له هوور څخه په خپل امريکايي چاپېريال کښې
  کمونست جوړوي۔ په دې جلا فرضي حالت کښې له خپل هېواد سره د
  هغه خيانت منو۔

  د دواړو مقدماتو له يوځاے کولو به دا نتيجه راووځي چې د
  هوور په روس کښې زېږېدل هم له خپل هېواد سره د هغه خيانت
  لازموي۔
\end{explain}
که هوور (په هماغه وخت چې په واقعيت کښې زېږېدلے و) د
متحده ايالاتو پر ځاے په روس کښې زېږېدلے وے، نو په شوروي
اتحاد کښې به لوے شوے وے او کمونست به شوے وے (دا فرض کړو)۔
د دې فرض له مخې لومړے مقدم رښتينے دے۔ په بېلګه کښې
دويم مقدم هم، چې په جلا توګه وکتل شي، رښتينے ګڼو۔ خو
د همدې فرضونو له مخې پايله نارښتينې ده: په دغه بديل حالت کښې
هوور که په شوروي اتحاد کښې زېږېدلے وے، نو کلک کمونست به وے
او (له خپل هېواد سره) غدار به نۀ وے۔ د سټالنکر--لويس شننه
د صدق د دغو ارزښتونو همدا وېش تاييدوي۔ زمونږ نړۍ ته هغه تر ټولو
نژدې ممکنه نړۍ چې يوازې د هوور د زېږون ځاے پکښې بدلېږي، هغه
ده چې هوور پکښې د شوروي اتحاد ښه وګړے لوېږي۔ دا هغه تر
ټولو نژدې ممکنه نړۍ ده چې د لومړي مقدم او پايلې مقدم پکښې
رښتينے دے؛ په هغې نړۍ کښې هوور د کمونست ګوند وفادار غړے
دے، نو غدار نۀ دے۔ خو د دويم مقدم د ارزولو لپاره بايد بله
نړۍ وګورو: تر ټولو نژدې نړۍ چې هوور پکښې کمونست دے، هغه
ده چې په متحده ايالاتو کښې زېږېدلے، بيا يې لورے بدل کړے او
په دې توګه غدار شوے دے۔\footnote{البته، د بېلګې د قوت د
  پوهېدو لپاره بايد ځينې مابعدالطبيعي او سياسي فرضونه ومنو؛
  د بېلګې په توګه، دا چې ممکنه وه هوور له روسي مور و پلار
  زېږېدلے وے، او دا سياسي فرض چې د نولس سوه پنځوسمو کلونو
  په متحده ايالاتو کښې کمونست کېدل خپل هېواد ته غداري وه۔}

\begin{prob}
  د واقعيت خلاف شرطونو د انتقاليت د ناکامۍ لپاره يوه
  قانع کوونکې، له لومړني تصور سره برابره بېلګه ومومئ۔
\end{prob}

\begin{ex}\ollabel{ex:trans-counterex}
  د کروي سيمو معنٰی پوهنه استنتاج نامعتبر ثابتوي، يعنې $p
  \cif q, q \cif r \Entails/ p \cif r$ لرو۔ مدل $\mModel{M}
  = \tuple{W, O, V}$ وګورئ چې $W = \{w, w_1, w_2\}$، $O_w = \{\{w\}, \{w,
  w_1\}, \{w, w_1, w_2\}\}$، $V(p) = \{w_2\}$، $V(q) = \{w_1, w_2\}$
  او $V(r) = \{w_1\}$ لري۔ د پاتې دوو نړۍو د سيمو نظام
  يوازې د هماغې نړۍ له يوغړي سټه جوړوو۔ يوه $p$-منونکې کروي سيمه $S = \{w, w_1,
  w_2\}$ شته او $p \lif q$ پکښې په ټولو نړۍو کښې رښتينې ده،
  نو $\mSat{M}{p
    \cif q}[w]$۔ يوه $q$-منونکې کروي سيمه $S' = \{w,
  w_1\}$ هم شته، او $q \lif r$ يې په ټولو نړۍو کښې رښتينے
  دے، نو $\mSat{M}{q \cif r}[w]$۔ خو يوازې
  بهرنۍ کروي سيمه د لومړي مقدم د رښتياوالي نړۍ لري۔ دغه
  $p$-منونکې کروي سيمه $\{w,
  w_1, w_2\}$ يوه نړۍ، يعنې~$w_2$، لري چې
  $\mSat/{M}{p \lif
    r}[w_2]$ پکښې برقرار دے۔ نو له لومړي مقدم څخه تر
  وروستۍ پايلې د واقعيت خلاف شرطي نارښتينې ده۔
\end{ex}

\begin{prob}
په \olref[cnt][min][tra]{ex:trans-counterex} کښې د ضدبېلګې
سره برابر د کروي سيمو انځور رسم کړئ۔
\end{prob}

\begin{prob}
  په \olref[cnt][min][tra]{ex:trans-counterex} کښې نړۍ $w_2$
  هغه ده چې هوور پکښې په روس کښې زېږېدلے، کمونست دے او
  غدار نۀ دے؛ او $w_1$ هغه نړۍ ده چې هوور پکښې په متحده
  ايالاتو کښې زېږېدلے، کمونست دے او غدار دے۔ په دې مدل کښې
  $w_1$ له~$w$ سره تر $w_2$ نژدې ده۔ آيا دا لازمي ده؟ داسې
  ضدبېلګه ورکولے شئ چې دا فرض نۀ کړي چې په روس کښې د هوور
  زېږېدل د هغه تر کمونست کېدو ډېر ليرې امکان دے؟ لارښوونه:
  د دواړو رښتينو شرطونو د شاهدې سيمې سره پرتله کړئ۔ که د
  دويم شرط په شاهده سيمه کښې د لومړي شرط مقدم رښتينے وي،
  د دغو دوو سيمو په تقاطع کښې به وروستۍ پايله څۀ ارزښت ولري؟
\end{prob}

\end{document}
